\documentclass[reqno]{exam}

\usepackage{amssymb, amsmath, amsthm, stmaryrd}
\usepackage{fullpage, parskip}
\usepackage{enumitem}
\usepackage{tikz}
\usepackage{ulem}
\usepackage{hyperref}

\title{Introduction to Computational Math Spring 2026 \\ Homework 08}
\date{Due: Friday, Mar 27, 2026, 11:59 PM}
\author{}

\begin{document}
\maketitle
\thispagestyle{empty}

Textbook = Driscoll and Braun (\url{https://fncbook.com/})

\begin{enumerate}

\item Textbook Exercise 5.4.6

\begin{solution}
Starting with the forward difference formula (5.4.2) for $f'$ at $x = 0$:
$$f'(0) \approx \frac{f(h) - f(0)}{h}$$

Applying this again to estimate $f''(0)$:
$$f''(0) \approx \frac{f'(h) - f'(0)}{h}$$

Now we use forward difference (5.4.2) for $f'(h)$:
$$f'(h) \approx \frac{f(2h) - f(h)}{h}$$

And backward difference (5.4.3) for $f'(0)$:
$$f'(0) \approx \frac{f(0) - f(-h)}{h}$$

Substituting these into the expression for $f''(0)$:
\begin{align*}
f''(0) &\approx \frac{1}{h}\left[\frac{f(2h) - f(h)}{h} - \frac{f(0) - f(-h)}{h}\right]\\
&= \frac{1}{h^2}\left[f(2h) - f(h) - f(0) + f(-h)\right]\\
&= \frac{f(-h) - f(0) - f(h) + f(2h)}{h^2}
\end{align*}

Therefore, the finite-difference formula is:
$$\boxed{f''(0) \approx \frac{f(-h) - f(0) - f(h) + f(2h)}{h^2}}$$

Or equivalently:
$$\boxed{f''(0) \approx \frac{f(2h) - f(h) - f(0) + f(-h)}{h^2}}$$
\end{solution}

\item Textbook Exercise 5.4.7
\begin{solution}
\textbf{(a)} The centered difference formula (5.4.5) is:
$$f'(x) \approx \frac{f(x+h) - f(x-h)}{2h}$$

We need to show that:
$$\lim_{h \to 0} \frac{f(x+h) - f(x-h)}{2h} = f'(x)$$

As $h \to 0$, both numerator and denominator approach 0, giving the indeterminate form $\frac{0}{0}$. Applying L'Hôpital's Rule (differentiating with respect to $h$):
$$\lim_{h \to 0} \frac{f(x+h) - f(x-h)}{2h} = \lim_{h \to 0} \frac{f'(x+h) \cdot 1 - f'(x-h) \cdot (-1)}{2}$$
$$= \lim_{h \to 0} \frac{f'(x+h) + f'(x-h)}{2} = \frac{f'(x) + f'(x)}{2} = f'(x)$$

Therefore, the centered formula converges to equality as $h \to 0$.

\textbf{(b)} Consider the general finite-difference formula:
$$f'(x) \approx \frac{1}{h}\sum_{k=-p}^{q} a_k f(x+kh)$$

For convergence as $h \to 0$, we need:
$$\lim_{h \to 0} \frac{\sum_{k=-p}^{q} a_k f(x+kh)}{h} = f'(x)$$

\textbf{Condition 1:} For the indeterminate form $\frac{0}{0}$ to occur (so L'Hôpital's Rule can be applied), the numerator must approach 0:
$$\lim_{h \to 0} \sum_{k=-p}^{q} a_k f(x+kh) = \sum_{k=-p}^{q} a_k f(x) = 0$$

This requires:
$$\boxed{\sum_{k=-p}^{q} a_k = 0}$$

\textbf{Condition 2:} Applying L'Hôpital's Rule:
$$\lim_{h \to 0} \frac{\sum_{k=-p}^{q} a_k f(x+kh)}{h} = \lim_{h \to 0} \frac{\sum_{k=-p}^{q} a_k \cdot k \cdot f'(x+kh)}{1} = f'(x) \sum_{k=-p}^{q} k a_k$$

For this to equal $f'(x)$:
$$\boxed{\sum_{k=-p}^{q} k a_k = 1}$$
\end{solution}

\item Textbook Exercise 5.5.2

\begin{solution}
The centered difference formula (5.4.3) is:
$$f'(x) \approx \frac{f(x+h) - f(x-h)}{2h}$$

The truncation error is:
$$\tau_f(h) = f'(x) - \frac{f(x+h) - f(x-h)}{2h}$$

Using Taylor series expansions about $x$:
\begin{align*}
f(x+h) &= f(x) + hf'(x) + \frac{h^2}{2}f''(x) + \frac{h^3}{6}f'''(x) + \frac{h^4}{24}f^{(4)}(x) + \frac{h^5}{120}f^{(5)}(x) + O(h^6)\\
f(x-h) &= f(x) - hf'(x) + \frac{h^2}{2}f''(x) - \frac{h^3}{6}f'''(x) + \frac{h^4}{24}f^{(4)}(x) - \frac{h^5}{120}f^{(5)}(x) + O(h^6)
\end{align*}

Subtracting:
$$f(x+h) - f(x-h) = 2hf'(x) + \frac{2h^3}{6}f'''(x) + \frac{2h^5}{120}f^{(5)}(x) + O(h^7)$$

Dividing by $2h$:
$$\frac{f(x+h) - f(x-h)}{2h} = f'(x) + \frac{h^2}{6}f'''(x) + \frac{h^4}{120}f^{(5)}(x) + O(h^6)$$

Therefore, the truncation error is:
$$\tau_f(h) = f'(x) - \left[f'(x) + \frac{h^2}{6}f'''(x) + \frac{h^4}{120}f^{(5)}(x) + O(h^6)\right]$$

$$\boxed{\tau_f(h) = -\frac{h^2}{6}f'''(x) - \frac{h^4}{120}f^{(5)}(x) + O(h^6)}$$

The first two nonzero terms are $-\frac{h^2}{6}f'''(x)$ and $-\frac{h^4}{120}f^{(5)}(x)$.
\end{solution}

\item Textbook Exercise 5.5.3
\begin{solution}
The second row of Table 5.4.2 gives the finite-difference formula:
$$f'(x) \approx \frac{-3f(x) + 4f(x+h) - f(x+2h)}{2h}$$

The truncation error is:
$$\tau_f(h) = f'(x) - \frac{-3f(x) + 4f(x+h) - f(x+2h)}{2h}$$

Using Taylor series expansions about $x$:
\begin{align*}
f(x) &= f(x)\\
f(x+h) &= f(x) + hf'(x) + \frac{h^2}{2}f''(x) + \frac{h^3}{6}f'''(x) + O(h^4)\\
f(x+2h) &= f(x) + 2hf'(x) + 2h^2f''(x) + \frac{4h^3}{3}f'''(x) + O(h^4)
\end{align*}

Computing the numerator:
\begin{align*}
&-3f(x) + 4f(x+h) - f(x+2h)\\
&= -3f(x) + 4\left[f(x) + hf'(x) + \frac{h^2}{2}f''(x) + \frac{h^3}{6}f'''(x)\right]\\
&\quad - \left[f(x) + 2hf'(x) + 2h^2f''(x) + \frac{4h^3}{3}f'''(x)\right] + O(h^4)\\
&= 2hf'(x) + (2h^2 - 2h^2)f''(x) + \left(\frac{2h^3}{3} - \frac{4h^3}{3}\right)f'''(x) + O(h^4)\\
&= 2hf'(x) - \frac{2h^3}{3}f'''(x) + O(h^4)
\end{align*}

Dividing by $2h$:
$$\frac{-3f(x) + 4f(x+h) - f(x+2h)}{2h} = f'(x) - \frac{h^2}{3}f'''(x) + O(h^3)$$

Therefore:
$$\boxed{\tau_f(h) = \frac{h^2}{3}f'''(x) + O(h^3)}$$

The first nonzero term is $\frac{h^2}{3}f'''(x)$.
\end{solution}

\item Textbook Exercise 5.5.5

\begin{solution}
The finite-difference formula (5.4.9) is:
$$f''(0) \approx \frac{f(-h) - 2f(0) + f(h)}{h^2}$$

The truncation error is:
$$\tau_f(h) = f''(0) - \frac{f(-h) - 2f(0) + f(h)}{h^2}$$

Using Taylor series expansions about $x = 0$:
\begin{align*}
f(h) &= f(0) + hf'(0) + \frac{h^2}{2}f''(0) + \frac{h^3}{6}f'''(0) + \frac{h^4}{24}f^{(4)}(0) + O(h^5)\\
f(-h) &= f(0) - hf'(0) + \frac{h^2}{2}f''(0) - \frac{h^3}{6}f'''(0) + \frac{h^4}{24}f^{(4)}(0) + O(h^5)
\end{align*}

Computing the numerator:
\begin{align*}
&f(-h) - 2f(0) + f(h)\\
&= \left[f(0) - hf'(0) + \frac{h^2}{2}f''(0) - \frac{h^3}{6}f'''(0) + \frac{h^4}{24}f^{(4)}(0)\right]\\
&\quad - 2f(0) + \left[f(0) + hf'(0) + \frac{h^2}{2}f''(0) + \frac{h^3}{6}f'''(0) + \frac{h^4}{24}f^{(4)}(0)\right] + O(h^5)\\
&= h^2f''(0) + \frac{h^4}{12}f^{(4)}(0) + O(h^5)
\end{align*}

Dividing by $h^2$:
$$\frac{f(-h) - 2f(0) + f(h)}{h^2} = f''(0) + \frac{h^2}{12}f^{(4)}(0) + O(h^3)$$

Therefore:
$$\tau_f(h) = -\frac{h^2}{12}f^{(4)}(0) + O(h^3)$$

Since the leading term of the truncation error is $O(h^2)$, the formula is \textbf{second-order accurate}.

$$\boxed{\tau_f(h) = O(h^2)}$$
\end{solution}

\item In this problem, we analyze a general-purpose algorithm called Richardson extrapolation that improves the convergence rate of finite difference approximations. 

Let $f$ be a smooth function whose derivative we wish to estimate at a point $a$. Suppose $D(h)$ is a finite difference approximation to $f'(a)$ with truncation error $O(h^p)$, where $p > 0$. 
Construct an improved estimate for $f'(a)$ by combining two finite difference approximations at different step sizes:
\[
R(h) = x D(h) + y D(2h),
\]
where $x$ and $y$ are constants to be determined.

\begin{itemize}
    \item Determine the conditions on $x$ and $y$ that ensure $R(h) \to f'(a)$ as $h \to 0$.
    \item Find specific values of $x$ and $y$ such that $R(h)$ converges to $f'(a)$ faster than $D(h)$ does. What is the order of convergence for $R(h)$?
    \item Suppose $D(h)$ is the forward difference approximation $D(h) = \frac{f(a+h) - f(a)}{h}$. Use Richardson extrapolation to derive an improved approximation for $f'(a)$ and determine its order of convergence.
\end{itemize}

\begin{solution}
    Since $D(h)$ has truncation error $O(h^p)$, we write:
\[
D(h) = f'(a) + c_1 h^p + c_2 h^{p+1} + O(h^{p+2}), \quad D(2h) = f'(a) + c_1 2^p h^p + c_2 2^{p+1} h^{p+1} + O(h^{p+2})
\]

\textbf{Part 1.} Substituting into $R(h) = x D(h) + y D(2h)$ gives $R(h) = (x + y)f'(a) + c_1(x + 2^p y)h^p + O(h^{p+1})$.
For convergence to $f'(a)$, we need $\boxed{x + y = 1}$.

\textbf{Part 2.} To eliminate the $O(h^p)$ term, we require $x + 2^p y = 0$. Solving the system gives:
\[
\boxed{x = \frac{2^p}{2^p - 1}, \quad y = \frac{-1}{2^p - 1}}
\]
With these values, $R(h) = f'(a) + O(h^{p+1})$, so the order of convergence is $\boxed{O(h^{p+1})}$.

\textbf{Part 3.} For forward difference, $p = 1$, so $x = 2$ and $y = -1$. Thus:
\[
R(h) = 2 \cdot \frac{f(a+h) - f(a)}{h} - \frac{f(a+2h) - f(a)}{2h} = \boxed{\frac{4f(a+h) - f(a+2h) - 3f(a)}{2h}}
\]
Order of convergence: $\boxed{O(h^2)}$.

\end{solution}

\item Jupyter notebook

\end{enumerate}


\section*{Quiz Information}

The quiz will be held on {\bf Tuesday, March 30, 2026} during the discussion section.

Quiz questions will be modifications of the homework problems and material discussed in class.

\textbf{Topics covered:} 

Chapters 5.4 and 5.5.

\begin{itemize}
    \item Finite difference approximations for derivatives
    \item Big O notation, error analysis and convergence rates of sequences
\end{itemize}




\end{document}